\documentclass[11pt]{article}

\usepackage[a4paper,margin=1in]{geometry}
\usepackage{amsmath}
\usepackage{amssymb}
\usepackage{amsthm}
\usepackage{bm}
\usepackage{booktabs}
\usepackage{array}
\usepackage{graphicx}
\usepackage{enumitem}
\usepackage{setspace}
\usepackage{hyperref}

\setstretch{1.15}

\title{
Latent Representation Learning for Financial Return Forecasting \\
\vspace{0.2cm}
\large A Temporal Variational Autoencoder and Latent Prediction Framework
}

\author{
Mushu
}

\date{\today}

\begin{document}

\maketitle

\tableofcontents

\newpage

% ============================================================
\section{Introduction}
% ============================================================

The objective of this project is to construct a probabilistic representation-learning pipeline for forecasting future company returns from historical financial information.

The central problem may be stated as follows.

Suppose a company possesses a sequence of historical monthly observations consisting of:
\[
x_1, x_2, \dots, x_t
\]
where each observation:
\[
x_i \in \mathbb{R}^d
\]
contains:
\begin{itemize}
    \item engineered financial factors,
    \item accounting ratios,
    \item profitability statistics,
    \item market variables,
    \item missingness indicators,
    \item and other numerical descriptors.
\end{itemize}

Given the historical sequence:
\[
X_{1:t}
=
(x_1,x_2,\dots,x_t)
\]
the task is to estimate:
\[
r_{t+1}
\]
where:
\[
r_{t+1}
\]
denotes the next-month gross return.

The project approaches this problem through latent representation learning.

Instead of directly learning:
\[
X_{1:t} \rightarrow r_{t+1}
\]
the model first learns a compressed probabilistic representation:
\[
X_{1:t} \rightarrow z
\]
and then learns:
\[
z \rightarrow r_{t+1}
\]

The pipeline therefore separates:
\begin{itemize}
    \item temporal representation learning,
    \item from downstream financial prediction.
\end{itemize}

This architecture allows:
\begin{itemize}
    \item dimensionality reduction,
    \item denoising,
    \item latent regime discovery,
    \item probabilistic encoding,
    \item and modular downstream experimentation.
\end{itemize}

% ============================================================
\section{Dataset Structure}
% ============================================================

The original dataset consists of monthly observations across multiple companies.

Each row corresponds to:
\[
(\text{company}, \text{month})
\]
and contains:
\begin{itemize}
    \item a timestamp,
    \item company identifier,
    \item numerical financial features,
    \item engineered indicators,
    \item and the target variable.
\end{itemize}

The target variable is:
\[
\texttt{monthly\_gross\_return}
\]
which represents the multiplicative gross return over the next month.

For example:
\[
1.05
\]
corresponds to a:
\[
5\%
\]
gain, while:
\[
0.95
\]
corresponds to a:
\[
5\%
\]
loss.

The dataset was partitioned company-wise into separate CSV files:
\[
\texttt{company\_data/company\_id.csv}
\]
where each file contains the complete historical trajectory of a single company.

Within each company:
\begin{itemize}
    \item rows are sorted chronologically,
    \item timestamps are converted to monthly datetime objects,
    \item and historical order is preserved.
\end{itemize}

% ============================================================
\section{Preprocessing Pipeline}
% ============================================================

\subsection{Sequence Construction}

The forecasting problem is transformed into a fixed-length sequence learning problem.

A sequence length:
\[
L = 12
\]
months is used.

For every company and every valid prediction month:
\[
t
\]
a historical sequence is constructed:
\[
X_t
=
(x_{t-L},\dots,x_{t-1})
\]

The associated prediction target is:
\[
y_t = r_t
\]

Thus each training example becomes:
\[
(X_t,y_t)
\]

The resulting tensor representation has the form:
\[
X
\in
\mathbb{R}^{N \times L \times d}
\]
where:
\begin{itemize}
    \item \(N\) is the number of sequences,
    \item \(L=12\) is sequence length,
    \item \(d\) is the feature dimension.
\end{itemize}

% ============================================================
\subsection{Handling Variable-Length Histories}

Some companies possess fewer than:
\[
12
\]
historical observations.

Instead of discarding these companies, padded sequences are constructed.

Suppose a company possesses only:
\[
k < L
\]
historical entries.

Then:
\[
L-k
\]
zero-vectors are prepended.

The final padded sequence therefore remains:
\[
L
\]
timesteps long.

A binary mask:
\[
m \in \{0,1\}^L
\]
is simultaneously constructed where:
\[
m_i
=
\begin{cases}
1 & \text{if timestep exists} \\
0 & \text{if timestep is padded}
\end{cases}
\]

This allows the reconstruction loss to ignore artificial padding.

% ============================================================
\subsection{Missing Features}

Some engineered financial features are missing for certain companies or months.

Rather than discarding incomplete observations:
\begin{itemize}
    \item missing values are filled,
    \item missingness indicators are preserved,
    \item and absent columns are inserted with zeros when necessary.
\end{itemize}

This preserves:
\begin{itemize}
    \item feature consistency,
    \item dimensional compatibility,
    \item and temporal continuity.
\end{itemize}

% ============================================================
\subsection{Chronological Splitting}

The dataset is partitioned chronologically:
\begin{itemize}
    \item training set,
    \item validation set,
    \item test set.
\end{itemize}

Chronological splitting is essential in financial forecasting because future information must never leak into the past.

Random shuffling across time would violate causal structure.

Thus:
\[
\text{train months}
<
\text{validation months}
<
\text{test months}
\]

% ============================================================
\section{Temporal Variational Autoencoder}
% ============================================================

\subsection{Motivation}

Financial sequences are:
\begin{itemize}
    \item high-dimensional,
    \item noisy,
    \item temporally dependent,
    \item and partially redundant.
\end{itemize}

Directly forecasting returns from raw features can therefore be unstable.

The Temporal Variational Autoencoder (TVAE) is used to learn a compressed latent representation:
\[
z
\]
that captures the underlying financial state of a company.

The model learns:
\[
X_{1:t}
\rightarrow
z
\]
where:
\[
z \in \mathbb{R}^k
\]
with:
\[
k \ll dL
\]

% ============================================================
\subsection{Encoder}

The encoder consists of a Gated Recurrent Unit (GRU).

Given:
\[
X
=
(x_1,\dots,x_L)
\]
the GRU processes the sequence recurrently:
\[
h_t
=
\text{GRU}(x_t,h_{t-1})
\]

The final hidden state:
\[
h_L
\]
summarizes the temporal dynamics of the entire sequence.

Two linear projections are then applied:
\[
\mu = W_\mu h_L + b_\mu
\]
\[
\log \sigma^2
=
W_\sigma h_L + b_\sigma
\]

These parameterize a Gaussian latent distribution:
\[
q(z|X)
=
\mathcal{N}(\mu,\sigma^2I)
\]

% ============================================================
\subsection{Latent Sampling}

The latent variable is sampled using the reparameterization trick:
\[
z
=
\mu
+
\sigma \odot \epsilon
\]
where:
\[
\epsilon
\sim
\mathcal{N}(0,I)
\]

This formulation allows gradients to propagate through stochastic sampling.

% ============================================================
\subsection{Decoder}

The decoder reconstructs the original sequence from:
\[
z
\]

The latent vector is projected into the GRU hidden space and repeated across timesteps.

The decoder therefore attempts to learn:
\[
z
\rightarrow
\hat X
\]

The reconstructed tensor:
\[
\hat X
\]
has identical dimensions to the original sequence tensor:
\[
X
\]

% ============================================================
\subsection{Loss Function}

The TVAE objective consists of two components.

\subsubsection{Reconstruction Loss}

The reconstruction term minimizes:
\[
\|X-\hat X\|^2
\]

Only valid timesteps contribute to loss:
\[
\mathcal{L}_{\text{recon}}
=
\frac{1}{|M|}
\sum
m_t
\|x_t-\hat x_t\|^2
\]

where:
\[
m_t
\]
is the timestep mask.

% ============================================================
\subsubsection{KL Divergence}

The latent distribution is regularized toward:
\[
\mathcal{N}(0,I)
\]

using:
\[
D_{KL}
\left(
q(z|X)
\|
p(z)
\right)
\]

which evaluates to:
\[
-\frac12
\sum
\left(
1+\log\sigma^2-\mu^2-\sigma^2
\right)
\]

% ============================================================
\subsubsection{Final Objective}

The complete objective becomes:
\[
\mathcal{L}
=
\mathcal{L}_{\text{recon}}
+
\beta
\mathcal{L}_{KL}
\]

where:
\[
\beta
\]
controls latent regularization strength.

% ============================================================
\section{Latent Representation Extraction}
% ============================================================

After TVAE training:
\begin{itemize}
    \item the decoder is discarded,
    \item and only the encoder is retained.
\end{itemize}

Each sequence:
\[
X_i
\]
is encoded into:
\[
\mu_i
\]

The latent dataset therefore becomes:
\[
Z
=
(\mu_1,\dots,\mu_N)
\]

where:
\[
Z
\in
\mathbb{R}^{N\times k}
\]

This latent space represents:
\begin{itemize}
    \item compressed financial states,
    \item temporal dynamics,
    \item and learned regime structure.
\end{itemize}

% ============================================================
\section{Latent Prediction Network}
% ============================================================

\subsection{MLP Predictor}

A Multi-Layer Perceptron (MLP) is trained on latent vectors.

The predictor learns:
\[
z
\rightarrow
\hat r_{t+1}
\]

The architecture consists of:
\begin{itemize}
    \item fully connected layers,
    \item ReLU nonlinearities,
    \item dropout regularization,
    \item and scalar regression output.
\end{itemize}

The network minimizes Mean Squared Error:
\[
\mathcal{L}_{\text{pred}}
=
\frac1N
\sum
(\hat r_i-r_i)^2
\]

% ============================================================
\section{Inference Pipeline}
% ============================================================

During inference:
\begin{enumerate}
    \item a company CSV file is loaded,
    \item rows are sorted chronologically,
    \item the most recent:
    \[
    12
    \]
    months are extracted,
    \item missing columns are inserted if necessary,
    \item the sequence is padded if required,
    \item the TVAE encoder computes:
    \[
    z
    \]
    \item the predictor computes:
    \[
    \hat r_{t+1}
    \]
\end{enumerate}

Thus the deployed forecasting system performs:
\[
X_{1:t}
\rightarrow
z
\rightarrow
\hat r_{t+1}
\]

% ============================================================
\section{Downstream Predictive Models}
% ============================================================

After latent representations are extracted from the Temporal Variational Autoencoder, downstream predictive models are trained to estimate future company returns.

The latent dataset consists of vectors:
\[
z_i \in \mathbb{R}^k
\]
where:
\[
k = 32
\]
in the current implementation.

Each:
\[
z_i
\]
represents a compressed latent financial state derived from:
\[
12
\]
months of historical company information.

The downstream learning problem therefore becomes:
\[
z_i
\rightarrow
r_{i+1}
\]

Two separate predictive architectures are explored:
\begin{enumerate}
    \item Multi-Layer Perceptron Predictor
    \item Latent Transformer Predictor
\end{enumerate}

Both models operate on the exact same latent representation space.

This allows the comparison to isolate:
\begin{itemize}
    \item the effect of downstream architecture,
    \item independent of representation learning.
\end{itemize}

% ============================================================
\subsection{Latent MLP Predictor}
% ============================================================

The first downstream predictor is a standard Multi-Layer Perceptron (MLP).

The objective of the MLP is:
\[
f_\theta :
\mathbb{R}^k
\rightarrow
\mathbb{R}
\]

where:
\[
f_\theta(z_i)
=
\hat r_{i+1}
\]

The network consists of:
\begin{itemize}
    \item fully connected linear layers,
    \item ReLU activation functions,
    \item dropout regularization,
    \item scalar regression output.
\end{itemize}

The architecture used in the implementation is:

\[
32
\rightarrow
128
\rightarrow
128
\rightarrow
1
\]

The forward computation may be written as:
\[
h_1
=
\phi(W_1 z + b_1)
\]

\[
h_2
=
\phi(W_2 h_1 + b_2)
\]

\[
\hat r
=
W_3 h_2 + b_3
\]

where:
\begin{itemize}
    \item \(\phi\) denotes the ReLU nonlinearity,
    \item \(h_1,h_2\) are hidden activations.
\end{itemize}

The MLP does not model temporal dynamics directly.

Instead:
\begin{itemize}
    \item temporal structure is assumed to already exist within the latent representation:
    \[
    z
    \]
    produced by the TVAE encoder.
\end{itemize}

Thus the MLP learns:
\[
z
\rightarrow
\hat r_{t+1}
\]

rather than:
\[
X_{1:t}
\rightarrow
\hat r_{t+1}
\]

This architecture serves as:
\begin{itemize}
    \item a baseline latent predictor,
    \item a latent-space regression model,
    \item and a representation-quality benchmark.
\end{itemize}

% ============================================================
\subsection{Latent Transformer Predictor}
% ============================================================

A second downstream architecture uses self-attention over latent representations.

Unlike traditional transformers operating on raw temporal sequences, the transformer in this project operates on:
\[
z
\]
directly.

The key observation is:
\begin{itemize}
    \item the Temporal VAE already models temporal dependencies using recurrent encoding,
    \item therefore the transformer is used to model latent interaction structure rather than temporal ordering.
\end{itemize}

% ============================================================
\subsubsection{Latent Tokenization}

The latent vector:
\[
z
\in
\mathbb{R}^{32}
\]
is reshaped into:
\[
8
\]
tokens of dimension:
\[
4
\]

Thus:
\[
z
\rightarrow
(z_1,z_2,\dots,z_8)
\]
where:
\[
z_i
\in
\mathbb{R}^4
\]

This produces a latent token sequence:
\[
Z
\in
\mathbb{R}^{8\times4}
\]

The intuition is that:
\begin{itemize}
    \item different subspaces of the latent representation may encode different financial factors,
    \item and self-attention may learn interactions between these latent components.
\end{itemize}

% ============================================================
\subsubsection{Token Projection}

Each latent token is projected into a higher-dimensional embedding space:
\[
e_i
=
W_p z_i + b_p
\]

where:
\[
e_i
\in
\mathbb{R}^{d_h}
\]

and:
\[
d_h = 64
\]
in the implementation.

Learnable positional embeddings are added:
\[
\tilde e_i
=
e_i + p_i
\]

even though the latent tokens are not temporal positions.

The positional embeddings instead provide:
\begin{itemize}
    \item latent token identity,
    \item latent subspace indexing,
    \item and structured attention ordering.
\end{itemize}

% ============================================================
\subsubsection{Self-Attention}

The transformer encoder applies multi-head self-attention:
\[
\text{Attention}(Q,K,V)
=
\text{softmax}
\left(
\frac{QK^\top}{\sqrt d}
\right)V
\]

where:
\[
Q,K,V
\]
are learned projections of latent token embeddings.

Self-attention allows the model to learn:
\begin{itemize}
    \item latent factor dependencies,
    \item nonlinear latent interactions,
    \item importance weighting between latent subspaces.
\end{itemize}

Multiple transformer encoder layers are stacked.

The implementation uses:
\begin{itemize}
    \item \(2\) attention heads,
    \item \(2\) encoder layers.
\end{itemize}

% ============================================================
\subsubsection{Pooling and Prediction}

After transformer encoding:
\[
H
=
(h_1,\dots,h_8)
\]

global average pooling is applied:
\[
\bar h
=
\frac18
\sum_{i=1}^8
h_i
\]

The pooled representation:
\[
\bar h
\]
is passed through a regression head:
\[
\bar h
\rightarrow
\hat r_{t+1}
\]

This produces the final return prediction.

% ============================================================
\subsection{Training Objective}
% ============================================================

Both downstream predictors minimize Mean Squared Error:
\[
\mathcal{L}_{\text{pred}}
=
\frac1N
\sum_{i=1}^N
(\hat r_i-r_i)^2
\]

Optimization is performed using Adam.

Validation loss is monitored across epochs and training logs are stored for reproducibility.

% ============================================================
\subsection{Model Comparison}
% ============================================================

The primary purpose of using both predictors is to compare:
\begin{itemize}
    \item simple latent regression,
    \item against attention-enhanced latent interaction modeling.
\end{itemize}

Since both models operate on identical latent representations:
\[
z
\]
performance differences can be attributed primarily to:
\begin{itemize}
    \item downstream predictive architecture,
    \item rather than representation learning itself.
\end{itemize}

This establishes a controlled experimental framework for evaluating:
\begin{itemize}
    \item latent-space predictive quality,
    \item attention usefulness,
    \item and representation structure.
\end{itemize}


% ============================================================
\section{Evaluation}
% ============================================================

Model performance is evaluated using:
\begin{itemize}
    \item Mean Squared Error,
    \item Mean Absolute Error,
    \item directional accuracy,
    \item correlation between predictions and targets.
\end{itemize}

Directional accuracy evaluates:
\[
\text{sign}(\hat r_t)
=
\text{sign}(r_t)
\]

which is financially important because:
\begin{itemize}
    \item predicting movement direction
    is often more useful than exact magnitude.
\end{itemize}

% ============================================================
\section{Future Work}
% ============================================================

Future extensions include:
\begin{itemize}
    \item latent transformers,
    \item hierarchical latent representations,
    \item volatility-aware objectives,
    \item probabilistic forecasting,
    \item attention-based latent interaction models,
    \item and reinforcement-learning trading policies.
\end{itemize}

The latent transformer extension specifically aims to model:
\[
z
\]
as a structured representation containing interacting latent financial factors.

% ============================================================
\section{Conclusion}
% ============================================================

This project constructs a complete representation-learning pipeline for financial return forecasting.

The framework:
\begin{itemize}
    \item preprocesses temporal company data,
    \item constructs chronological sequence tensors,
    \item learns probabilistic latent financial states,
    \item and performs downstream return prediction.
\end{itemize}

The architecture separates:
\begin{itemize}
    \item temporal representation learning,
    \item from predictive modeling.
\end{itemize}

This modularity enables:
\begin{itemize}
    \item systematic experimentation,
    \item latent-space analysis,
    \item and downstream model comparison.
\end{itemize}

\end{document}